% Transcription — fonds Alexandre Grothendieck, shelfmark 144, batch 6 (pages 101-120).
% Established from the facsimile archives/batches/144/batch-06.pdf, per the skill
% transcribe-grothendieck.
%
% Pass: Opus 5.5 (claude-opus-5-5), 2026-10-03 — first pass, unchecked against the pages by a human.
%
% Transcribed: 102, 104, 106, 108, 110, 112, 114 (his numbered sheets 11-16 and
% one unnumbered sheet), and 117-120 (two letters of P. Deligne — third-party
% correspondence, summarised in notes and not transcribed (RIGHTS.md) — the first dated
% 27/3/81). Skipped: 101 (blank), 103, 113, 115 (a student's copy on the
% icosahedron, dated 16 June 1978, on the backs of his sheets), 105, 107, 109,
% 111 (USTL administrative notes, 105 with only show-through), 116 (a cover:
% its inscription becomes the section heading).
%
% \page{} numbers the position in the folder. From position 117 the pencil
% numbers run one ahead (117 is pencilled 118, 120 is pencilled 121); a note
% at page 117 records it.
\documentclass[11pt,a4paper]{article}
\input{../preamble/grothendieck}

\folder{144}
\batch{6}
\pages{101}{120}
\foldertitle{[Suite autour de Teichmüller dont] Teichmülleries, 1981-1982 : notes manuscrites (1981-1983, s.d.), lettres (1981, s.d.).}
\dating{1981-1983}
\watermark{Édition de démonstration}

\begin{document}

\note{Les feuillets 104 à 114 portent sa numérotation 11 à 16 : le calcul a commencé avant ce lot (feuillets 1 à 10 de sa main). Il porte sur le groupe $\Pi_{0,3}$ de $\mathbb{P}^1$ privé de $0, 1, \infty$, ses lacets $l_0, l_1, l_\infty$ (avec $l_\infty l_1 l_0 = 1$), le revêtement de degré 2 donné par $z \mapsto (2z-1)^{2}$ et les relèvements $\hat{\varphi}$ au complété profini.}

\page{102}
\drawing{Une sphère (la droite projective complexe) : un grand cercle horizontal portant les points $0$, $1/2$, $1$, $2$, $-1$, le point $\infty$ marqué sur la face arrière ; un méridien passant par $1/2$ ; un fuseau hachuré entre deux méridiens.}

\struck{$e_0$ — $e_1$}
\[
\begin{array}{ll}
e_0 \longmapsto e_0 & \qquad e_0 \longmapsto e_1 \\
e_1 \longmapsto e_1 - e_0 & \qquad e_1 \longmapsto -e_0
\end{array}
\]
\struck{$B$} \note{lettre biffée, lecture douteuse}
\[
e_0 \longmapsto e_0 + e_1, \qquad e_1 \longmapsto -e_0
\]
\[
\begin{pmatrix} 1 & \lambda \\ 0 & 1 \end{pmatrix}
\qquad
\begin{pmatrix} \zeta & 0 \\ 0 & \zeta^{*} \end{pmatrix}
\]
\note{dans la seconde matrice, un $\lambda$ est surchargé par $\zeta$ ; l'exposant de $\zeta^{*}$ est \uncertain{$*$}.}

\struck{\ill{}}
\[
\det(u - \lambda\,\mathrm{id}) = \det \begin{pmatrix} 1-\lambda & -1 \\ 1 & -\lambda \end{pmatrix} = \lambda^{2} - \lambda + 1
\]
racines $\mapsto$ racines primitives 6\supplied{ièmes} de $1$

\note{trait horizontal sur toute la largeur de la page.}

$l_0 \quad l_{1/2} \quad l_1 \quad l_\infty \;\big/\; \dot{\sigma}_\infty,\ \bar{\tau},\ \sigma_\infty \tau$ \note{les accents sur $\sigma_\infty$ et $\tau$ sont lus \uncertain{point} et \uncertain{barre}.}

\[
\boxed{\,l_\infty\, l_1\, l_{1/2}\, l_0 = 1\,}
\]

\[
z \longmapsto 1 - z
\]
\[
\begin{array}{rcl}
g : \mathbb{P}^1 & \longrightarrow & \mathbb{P}^1 \\
\infty & \longmapsto & \infty \\
1/2 & \longmapsto & 0 \\
0,\ 1 & \longmapsto & 1
\end{array}
\qquad
\begin{array}{l}
g(z) = (2z-1)^{2} \\
g(1-z) = (1-2z)^{2} \\
\end{array}
\]
$(2z - 1\text{\struck{\ill{}}})^{2}$ \note{le dernier terme est surchargé, lecture incertaine.}
\[
\begin{array}{l}
l_0 \longmapsto -1 \\
l_1 \longmapsto +1 \\
l_\infty \longmapsto -1
\end{array}
\]

\page{104}
\note{numéro de sa main en tête de feuillet : 11.}
$\lambda'_0 \quad \lambda'_{1/2} \quad \lambda'_1 \quad \lambda'_\infty$
\[
\pi' \subset \widehat{\Pi}_{0,3}, \qquad \pi' \longrightarrow \widehat{\Pi}_{0,3}
\]
\note{le second $\widehat{\Pi}_{0,3}$ est au bout d'une flèche oblique issue de $\pi'$.}

\drawing{Un point base d'où partent des chemins vers les points $0$, $1/2$, $1$ d'une droite horizontale, chacun contournant le point par un petit lacet fléché ; plus loin sur la droite un point $z$ et, au bout, $\infty$.}

\[
(-2\bar{j} - 1)^{2} = 4\bar{j}^{2} + \ldots
\]
\[
\lambda^{k}_{1/2} \text{ \struck{\ill{}} } \longmapsto l_0
\qquad
\text{\struck{$2z-1$}} \longmapsto \lambda
\]

\drawing{Deux sphères. Sur la première, un cercle horizontal portant $-1$, $0$, $1$, $\infty$ ; lacets fléchés $\lambda'_{1/2}$, $\lambda'_1$, $\lambda'_\infty$ autour des points, et un chemin marqué $3$. Sur la seconde, le pôle marqué $-\bar{j}$, un méridien épais descendant au point $1/2$ du cercle équatorial (en tirets), les points $0$, $1$, $\infty$ ; une petite flèche tourne autour de $1/2$. À droite, un fuseau (lentille) coupé par deux segments perpendiculaires.}

\[
\text{\struck{$2z-1$}} \qquad (2x - 1 + 2iy)^{2} = (2x-1)^{2} - 4y^{2} + 4iy(2x-1)
\]
\[
y = 0 \qquad x = 1/2
\]
\[
\text{\struck{\ill{}$(z) = 1 - z$}} \quad (-2\bar{j} + 1)^{2} = 4\bar{j}^{2} + 1 - 4\bar{j} = 4(\bar{j}^{2} - \bar{j})
\]
\struck{\ill{}}
\[
(2\bar{j} - 1)^{2} = 4(\bar{j}^{2} - \bar{j}) + 1 = -3
\]
\[
\bar{j}^{2} - \bar{j} + 1 = (\bar{j} - \alpha)(\bar{j} - \bar{\alpha})
\]
\note{la dernière égalité est lue avec doute ($\bar{\alpha}$ ou $\beta$) ; en dessous, un fragment biffé : \struck{$\alpha \ldots$}.}

\[
\begin{array}{lll}
\lambda_0 & \longmapsto & \text{\struck{$(l_\infty^{2}\, l_1\, l_0^{2})^{-1}$}} = l_0^{-2}\, l_1^{-1}\, l_\infty^{-2} \\
\lambda_{1/2} & \longmapsto & l_0^{2} \\
\lambda_1 & \longmapsto & l_1 \\
\lambda_\infty & \longmapsto & l_\infty^{2}
\end{array}
\]
\[
\begin{aligned}
&= l_0^{-1}\, l_\infty^{-1} = l_0^{-1}\, l_1\, l_0 \\
&= (l_\infty l_0)^{-1} \\
&\text{\struck{$= \bigl(\mathrm{int}(l_0^{-1})(l_0 l_\infty)\bigr)^{-1}$}} \\
&= \mathrm{int}(l_0^{-1})\bigl((l_0 l_\infty)^{-1}\bigr) \\
&= \mathrm{int}(l_0^{-1})(l_1)
\end{aligned}
\]
\note{il note \uncertain{$\mathrm{int}$} l'automorphisme intérieur ; sous $(l_0 l_\infty)^{-1}$ il écrit $l_1$.}
\[
\text{\struck{$= l_0^{-1} l_1 l_0$}} \qquad
l_1\, l_0\, l_\infty = 1, \qquad l_1^{-1} = l_0\, l_\infty, \qquad
\text{\struck{$l_1 = l_\infty^{-1}\, l_0^{-1}$}}
\]
\[
l_0^{-1}(l_1 l_\infty \quad l_0 \qquad
\text{\struck{$ab\,(a^{-1}b^{-1})\,(ba)\,ab$}}
\]

\page{106}
\note{numéro de sa main : 12.}
\[
\pi' \subset \Pi_{0,3}
\]
\uncertain{engendré par}
\[
\begin{array}{l}
\text{\struck{\ill{}}} \\
\lambda_0 = \mathrm{int}(l_0^{-1})\, l_1 \\
\lambda_{1/2} = l_0^{2} \\
\lambda_1 = l_1 \\
\lambda_\infty = l_\infty^{2}
\end{array}
\qquad
\begin{array}{l}
\Pi_{0,3} \longrightarrow \{\pm 1\} \longrightarrow 1 \\
l_0,\ l_\infty \longmapsto -1 \\
l_1 \longmapsto +1
\end{array}
\]

\begin{tikzcd}[column sep=small, row sep=small, nodes={font=\scriptsize}]
1 \arrow[r] & \pi' \arrow[r, hook] \arrow[d, "\text{épi}"] & \Pi_{0,3} \arrow[r] \arrow[d] & \{\pm 1\} \arrow[r] \arrow[d] & 1 \\
1 \arrow[r] & \Pi_{0,3} \arrow[r] & \Pi_{0,3}\cdot\{1,\sigma_\infty\} \arrow[r] & \{\pm 1\} \arrow[r] & 1
\end{tikzcd}

\note{les $1$ de gauche et de droite de la seconde ligne sont à demi coupés sur la feuille.}

$l_0 \longmapsto \sigma_\infty$ ?
\[
\begin{array}{lll}
\lambda_0 = \mathrm{int}(l_0^{-1})\, l_1 = l_0^{-1} l_1 l_0 & \longmapsto & l_0 \\
\lambda_{1/2} = l_0^{2} & \longrightarrow & 1 \\
\lambda_1 = l_1 & \longrightarrow & l_1 \\
\lambda_\infty = l_\infty^{2} & \longrightarrow & l_\infty
\end{array}
\]
\note{dans les membres de gauche, des indices sont surchargés ($l_0$ sur $l_1$, $l_1$ sur $l_0$…) ; on donne la forme finale.}

\struck{$l_\infty^{2}\, l_1\, l_0\, l_0^{-1} l_1 l_0$} \note{sous l'expression biffée, il écrit $l'_\infty$.} \quad $\hat{\varphi}$ :
\[
\boxed{
\begin{array}{l}
l_0 \longmapsto \sigma_\infty \\
l_1 \longmapsto l_1 \\
l_\infty \longmapsto \varepsilon_\infty
\end{array}}
\]
\note{deux points d'interrogation en marge de l'encadré, en face de $\sigma_\infty$ et de $\varepsilon_\infty$.}
\[
\varepsilon_\infty\, l_1\, \sigma_\infty = \varepsilon_\infty\, \varepsilon_1^{2}\, \sigma_\infty
\]
\[
\begin{array}{l}
\varepsilon_0 = \sigma_\infty \rho \\
\varepsilon_1 = \rho\, \sigma_\infty \\
\varepsilon_\infty = \rho^{-1} \sigma_\infty \rho^{-1}
\end{array}
\]
\struck{$\rho^{-1} \sigma_\infty \rho^{-1}\, \rho\, \sigma_\infty\, \rho\, \sigma_\infty\, \sigma_\infty$}

\page{108}
\note{numéro de sa main : 13.}
\[
\widehat{\Pi}_{0,3} \xrightarrow{\ \hat{\varphi}_{1/2}\ } \widehat{\Pi}_{0,3}\cdot\{1,\sigma_\infty\}
\]
\[
\hat{\varphi}\bigl(u(x)\bigr) = u\,\hat{\varphi}(x)
\]
\note{au-dessus de $u$, relié par un trait : $\mathrm{int}(\varphi(g))$, lecture \uncertain{$\varphi(g)$} pour l'argument.}

1) $x = l_0$
\[
\begin{aligned}
u(x) &= \mathrm{int}(g_0)\, l_0^{p} \\
\hat{\varphi}\bigl(u(x)\bigr) &= \mathrm{int}\bigl(\hat{\varphi}(g_0)\bigr)\, \sigma_\infty^{p} \\
\hat{\varphi}(x) &= \sigma_\infty \qquad u\bigl(\hat{\varphi}(x)\bigr) = \sigma_\infty
\end{aligned}
\]
\[
\bigl(\mathrm{int}(\hat{\varphi}(g_0))\bigr)(\sigma_\infty^{p}) = \text{\struck{\ill{}}}\,(\sigma_\infty)
\]
\note{sous $\sigma_\infty^{p}$, un signe d'égalité vertical renvoie à $\sigma_\infty$.}

\marginal{$\hat{\varphi}(g_0)$ \uncertain{commute à} $\sigma_\infty$ ; $\alpha$ \uncertain{dépend du choix de $g_0$ modulo} \ill{}.}

\[
\text{\struck{\ill{}}}\qquad
\boxed{\,\hat{\varphi}_{1/2}(g_0) = \sigma_\infty^{\alpha}\,}
\]

2) $x = l_1$
\[
\begin{aligned}
\hat{\varphi}\bigl(u(x)\bigr) &= \hat{\varphi}\bigl(\mathrm{int}(g_1)\, l_1^{p}\bigr) = \mathrm{int}\bigl(\hat{\varphi}(g_1)\bigr)(l_1^{p}) \\
u\bigl(\hat{\varphi}(x)\bigr) &= u(l_1) = \mathrm{int}(g_1)(l_1^{p})
\end{aligned}
\]
\[
\text{\struck{$\hat{\varphi}_{1/2}(g_1) = g_1\, l_1^{\beta}$}}
\]
\note{cette ligne encadrée est barrée d'un trait horizontal, avec une accolade ondulée dessous.}
\[
g_1 = g\,\rho(g_0) = \sigma_\infty(g_0\, l_0^{\gamma}) \;=\; \sigma_\infty(g_0)\, l_1^{\gamma}
\]
\note{la lecture de $g\,\rho(g_0)$ est douteuse ; dans la seconde égalité il écrit $l_1$ là où la première porte $l_0$.}
\[
\hat{\varphi}\bigl(\sigma_\infty(g_0)\, l_1^{\gamma}\bigr) = \sigma_\infty(g_0)\, l_1^{\beta+\gamma}
\]
\[
\boxed{\,\hat{\varphi}_{1/2}\bigl(\sigma_\infty(g_0)\bigr) = \sigma_\infty(g_0)\, l_1^{\beta}\,}
\]
\[
\sigma_\infty(g_0)^{-1}\, g_0 = h_0, \qquad \sigma_\infty(g_0) = g_0\, h_0^{-1}
\]
\marginal{$\beta$ \uncertain{ne dépend que du choix de $g_0$} \uncertain{modulo} $L_0$.}
\[
\begin{aligned}
\hat{\varphi}\bigl(g_0\, h_0^{-1}\, l_1^{\gamma}\bigr) &= \text{\struck{\ill{}}}\; g_0\, h_0^{-1}\, l_1^{\beta+\gamma} \\
\hat{\varphi}(g_0)\, \hat{\varphi}(h_0^{-1})\, l_1^{\gamma} &= g_0\, h_0^{-1}\, l_1^{\beta+\gamma}
\end{aligned}
\]
\note{sous $\hat{\varphi}(g_0)$, une accolade renvoie à $\sigma_\infty^{\alpha}$. Les deux dernières lignes sont barrées en partie d'un trait oblique.}
\[
\hat{\varphi}(h_0) = l_1^{-\beta}\, h_0\, g_0^{-1}\, \sigma_\infty^{\alpha}
\]

\[
\begin{aligned}
&\text{3) } x = l_\infty \qquad
\hat{\varphi}\bigl(u(x)\bigr) = \hat{\varphi}\bigl(\mathrm{int}(g_\infty)\, l_\infty^{p}\bigr) = \mathrm{int}\bigl(\hat{\varphi}(g_\infty)\bigr)\, \varepsilon_\infty^{p} \\
&\phantom{\text{3) } x = l_\infty \qquad} u\,\hat{\varphi}(x) = u(\varepsilon_\infty) = \mathrm{int}(g_\infty)\, \varepsilon_\infty^{p}
\end{aligned}
\]
\[
\text{\struck{\ill{}}}\;\hat{\varphi}(g_\infty) = g_\infty
\]
\note{le cas 3) est entouré à gauche d'un crochet.}

\page{110}
\note{numéro de sa main : 14.}
\[
\pi' \longrightarrow \pm 1, \qquad \lambda_0,\ \lambda_{1/2},\ \lambda_1,\ \lambda_\infty \longmapsto -1
\]
\[
\Pi_{0,3} \longrightarrow \mathbb{Z}/4\mathbb{Z}
\]
\[
\left.
\begin{array}{l}
\text{\struck{$\lambda$}}\; l_0^{-1} l_1 l_0 \\
l_0^{2} \\
l_1 \\
l_\infty^{2}
\end{array}
\right\}
\longmapsto 2 \bmod 4
\qquad\qquad
\begin{array}{l}
l_0 \longmapsto \alpha_0 \\
l_1 \longmapsto \alpha_1 = 2 \bmod 4 \\
l_\infty \longmapsto \alpha_\infty
\end{array}
\]
\note{à gauche, un trait vertical réunit les quatre éléments et une seule flèche part vers $2 \bmod 4$.}
\[
\begin{array}{ll}
\alpha_0 + \alpha_\infty \equiv 2 & (4) \\
2\alpha_0 \equiv 2 & (4) \\
2\alpha_\infty \equiv 2 & (4) \\
\alpha_0 \equiv \pm 1 & (4) \\
\alpha_\infty \equiv \pm 1 & (4)
\end{array}
\qquad
\begin{array}{l}
\text{sps } \alpha_0 \equiv 1 \\
\text{alors } \alpha_\infty \equiv 1
\end{array}
\]
\note{« sps » : sans doute « supposons ».}
\[
\boxed{
\begin{array}{l}
l_0,\ l_\infty \longmapsto 1 \bmod 4 \\
l_1 \longmapsto 2 \bmod 4
\end{array}}
\]

\begin{tikzcd}[column sep=small, row sep=small, nodes={font=\scriptsize}]
1 \arrow[r] & \pi'' \arrow[r] \arrow[dr] & \Pi_{0,3} \arrow[r] & \mathbb{Z}/4\mathbb{Z} \arrow[r] & 1 \\
 & & \pi''/\Delta(\lambda_0^{2},\lambda_1^{2},\lambda_\infty^{2}) & &
\end{tikzcd}

\note{devant le quotient, un début biffé : \struck{$\lambda_0 \cdot \pi''_{\ldots}$}. Sous les trois générateurs du quotient, des signes d'égalité verticaux renvoient à $l_0^{-1} l_1^{2} l_0$, $l_1^{2}$, $l_\infty^{4}$ ; la lecture de $\Delta$ est \uncertain{douteuse}.}

\begin{tikzcd}[column sep=small, row sep=small, nodes={font=\scriptsize}]
 & \mathrm{Autext}_{\mathrm{car}}(\pi'')' & & & \\
1 \arrow[r] & \widehat{\mathcal{S}}^{+}_{0,3} \arrow[r] \arrow[u] & M_{0,3} \arrow[r] \arrow[d] & \Pi \arrow[r] & 1 \\
 & & N'_{1,1} & &
\end{tikzcd}

\note{l'indice de $\mathrm{Autext}$ est lu \uncertain{car} ; la lettre $\widehat{\mathcal{S}}$ (une S cursive coiffée d'un chapeau, au bout de la flèche verticale) est d'une lecture \uncertain{incertaine}. Sur la flèche $M_{0,3} \to \Pi$, une étiquette biffée : \struck{$l_0$}.}

\page{112}
\note{numéro de sa main, cerclé : 15.}
\[
\begin{array}{lll}
l_0^{-1} l_1 l_0 = \lambda_0 & \longmapsto & 1 \\
l_0^{2} = \lambda_{1/2} & \longrightarrow & l_0 \\
l_1 = \lambda_1 & \longmapsto & l_1 \\
l_\infty^{2} = \lambda_\infty & \longmapsto & l_\infty
\end{array}
\qquad\qquad
\begin{array}{lll}
l_0 & \longmapsto & \varepsilon_0 \\
l_1 & \longmapsto & l_1 \\
l_\infty & \longmapsto & \varepsilon_\infty
\end{array}
\]
\[
\varepsilon_\infty\, l_1\, \varepsilon_0 = 1 \quad \text{i.e.} \quad
\underbrace{\rho^{-1}\sigma_\infty\rho^{-1}}\;\underbrace{\rho\,\sigma_\infty\,\rho\,\sigma_\infty}\;\sigma_\infty\,\rho = \rho^{-1}\sigma_\infty
\]
\note{dans le produit, les facteurs qui se simplifient sont barrés un à un de traits obliques ; les accolades sont les siennes.}
\[
\begin{array}{l}
\varepsilon_0 = \sigma_\infty\, \rho \\
\varepsilon_1 = \rho\, \sigma_\infty \\
\varepsilon_\infty = \rho^{-1} \sigma_\infty\, \rho^{-1}
\end{array}
\]

\page{114}
\note{numéro de sa main : 16.}
\[
\begin{array}{lll}
\lambda_0 = l_0^{-1} l_1 l_0 & \longmapsto & l_1 \\
\lambda_{1/2} = l_0^{2} & \longmapsto & l_\infty \\
\lambda_1 = l_1 & \longmapsto & l_0 \\
\lambda_\infty = l_\infty^{2} & \longmapsto & 1
\end{array}
\qquad\qquad
\begin{array}{lll}
l_0 & \longmapsto & \varepsilon_\infty \\
l_1 & \longmapsto & l_0 \\
l_\infty^{\ast} & \longrightarrow & \sigma_\infty
\end{array}
\]
\note{les images $l_1$ et $1$ de la première colonne sont écrites par-dessus d'autres signes ; l'exposant de $l_\infty^{\ast}$ est un trait \uncertain{(peut-être $\infty$)}.}
\[
\begin{aligned}
(l_1 \varepsilon_0)^{-1} &= (\rho\,\sigma_\infty\, \rho^{-1}\ \ldots)^{-1} = \sigma_\infty\, \rho^{-1} \\
(l_0\, \varepsilon_\infty)^{-1} &= (\sigma_\infty\rho\,\sigma_\infty\rho\,\rho^{-1}\sigma_\infty\rho^{-1}) = \sigma_\infty
\end{aligned}
\]
\note{ici encore les facteurs simplifiés sont barrés de traits obliques ; dans la première ligne, le facteur barré qui suit $\sigma_\infty$ est \uncertain{$\rho^{-1}\sigma_\infty\sigma_\infty\rho$}.}
\[
\hat{\varphi}_\infty : \widehat{\Pi}_{0,3} \longrightarrow \widehat{\Pi}_{0,3}\cdot\{1,\sigma_\infty\}
\]

\[
\mathrm{Sl}(2,\mathbb{Z}/12)' \simeq \bigl(\mathrm{Sl}(2,\mathbb{Z}/4) \times \mathrm{Sl}(2,\mathbb{Z}/3)\bigr)/\pm 1
\]

\begin{tikzcd}[column sep=small, row sep=small, nodes={font=\scriptsize}]
\mathrm{Sl}(2,\mathbb{Z}/12)' \arrow[d] \\
\mathfrak{S}_3 \simeq \mathrm{Sl}(2,\mathbb{Z}/2)
\end{tikzcd}

\note{au-dessus, isolé, $\hat{\varphi}_\infty$. Le « $12$ » du premier membre est lu \uncertain{$\mathbb{Z}/12$} : la page porte « $(2,12)'$ ».}

\begin{tikzcd}[column sep=small, row sep=small, nodes={font=\scriptsize}]
\widehat{\mathfrak{S}}^{+}_{0,3} \arrow[d] & \\
\mathrm{Sl}(2,\widehat{\mathbb{Z}})' \arrow[r] & \mathrm{Sl}(2,\mathbb{Z}/2\mathbb{Z}) \arrow[r] & 1
\end{tikzcd}

\note{sous $\mathrm{Sl}(2,\widehat{\mathbb{Z}})'$, un signe $\wr$ renvoie à un second $\widehat{\mathfrak{S}}^{+}_{0,3}$, surchargé.}

\struck{$(c,d)$ \quad
$c \equiv 0\ (3)$ : $D = d^{2} \in \{0,1\}$, $0$ exclus ;
$c \equiv 1\ (3)$ : $D = 1 + d^{2} + d \in \{1,0\}$, $0$ exclus ;
$c \equiv -1\ (3)$ : $D = 1 + d^{2} - d \in \{1,0\}$, $0$ exclus !!!}
\note{tout le bloc est biffé d'un grand zigzag ; dans la deuxième ligne, un membre intermédiaire est surchargé.}

\section{Calculs divers sur Teichmüller — Lettre Deligne sur groupes profinis}
\note{Titre de sa main, au crayon, sur un feuillet par ailleurs blanc (p. 116 de la cote) : « Calculs divers sur Teichmüller » et, en dessous, « Lettre Deligne sur groupes profinis ». Le feuillet fait couverture ; la première ligne peut désigner aussi les calculs qui précèdent. Les pages suivantes sont deux lettres de P.~Deligne à Grothendieck (« Schurik »), qui ne sont pas transcrites.}

\page{117}
\note{À partir d'ici, le crayon de l'archiviste porte un numéro d'avance sur la position dans la cote (ce feuillet porte 118, le dernier du lot 121). Les numéros employés pour \texttt{page} sont ceux de la position, que suit le fac-similé.}

\note{Lettre manuscrite de Pierre Deligne à Grothendieck, datée « 27/3/81 », pp.~117--118, signée « Pierre » : pro-groupes libres (renvoi au ch.~I de Serre, \emph{Cohomologie galoisienne}), éléments d'ordre fini et éléments centraux d'un complété profini de groupe libre, centre des sous-groupes d'indice fini de $\mathrm{Gal}(\overline{\mathbb{Q}}/\mathbb{Q})$, sous-groupes finis du groupe de Teichmüller (référence à Kerckhoff, Bull.\ AMS) ; un paragraphe personnel en russe. Correspondance d'un tiers, non transcrite (RIGHTS.md).}

\page{118}

\page{119}
\note{Seconde lettre de Deligne, sans date, pp.~119--120 et suite au lot~7 (pp.~121--123) : un résumé de la théorie du corps de classes, locale, globale et local-global (renvoi à Tate, « Number theoretic background », Corvallis), puis le théorème de Shimura--Taniyama sur la multiplication complexe. Correspondance d'un tiers, non transcrite (RIGHTS.md).}

\page{120}

\end{document}
